% chapters/ch03-keys-identity.tex
\chapter{Keys and Identity}

\begin{learningobjectives}
In this chapter you will learn how identity drives relational design.
You will be able to choose appropriate keys, distinguish natural and surrogate identity, and design composite keys for association relations.
You will also learn how key choices affect correctness, not just convenience.
\end{learningobjectives}

\section{Why identity comes first}
Identity is the backbone of a relational model.
If you cannot identify an entity instance, you cannot reference it, and constraints cannot bind it to anything else.
All relationships depend on identity, and all constraints depend on references.

In practice, unclear identity produces duplicate rows that look different but represent the same thing.
Once that happens, every rule is weakened.
Two rows for the same patient mean two appointment histories, two sets of constraints, and two conflicting truths.

\begin{principlebox}
Identity is not optional.
Every entity must have a stable way to be referenced.
\end{principlebox}

\section{Natural keys and surrogate keys}
A \term{natural key} is a real-world identifier that already exists in the domain.
An ISBN can identify a book title.
A national health number might identify a patient.
These keys can be meaningful, but they are not always stable or available.

A \term{surrogate key} is an identifier created for the system (often a numeric id).
It carries no meaning, but it is stable and easy to reference.
Many systems use surrogate keys for consistency and then enforce natural uniqueness separately.

\begin{examplebox}
\textbf{Library title identity}\par\smallskip
If ISBN is unique and stable in your domain, it can be a natural key:\par
\[
\rel{Title}(\underline{isbn}, title, author)
\]
If ISBN is missing or unreliable, use a surrogate key and enforce uniqueness on ISBN when present.
\end{examplebox}

\begin{examplebox}
\textbf{Clinic patient identity}\\
Patients may not have a stable external identifier at first contact.
A surrogate key avoids ambiguity:
\[
\rel{Patient}(\underline{patient\_id}, name, birth\_date)
\]
Other identifiers can be stored as attributes with uniqueness constraints if required.
\end{examplebox}

\section{Key stability and lifecycle}
Keys live through events: creation, correction, merges, and splits.
If a key is not available at creation time, a surrogate key is usually safer.
If a key can change, you must plan for updates across every reference.

\subsection{Creation time and missing identifiers}
Some identifiers arrive later.
Clinic patients may receive a national identifier after the first visit.
Library copies may receive a barcode after cataloguing.
If you model these identifiers as primary keys, you will be forced to invent placeholders or delay record creation.

\subsection{Merges and splits}
Real entities can merge or split.
Two patient records can be merged when a duplicate is discovered.
A university course code can be split into two courses when curriculum changes.
These events are manageable when identity is stable and references are consistent.
They are chaotic when identity is tied to mutable attributes.

\begin{examplebox}
\textbf{Merging duplicate patients}\par\smallskip
If two rows represent the same person, you need to preserve references:
\[
\rel{Patient}(\underline{patient\_id}, name, birth\_date, national\_id? \;[\text{unique}])
\]
Merging uses a single \texttt{patient\_id} while keeping optional external identifiers for matching.
\end{examplebox}

\section{Key quality criteria}
Good keys share a few properties.
These criteria are not formal rules, but they prevent most identity failures.

\begin{principlebox}
\textbf{Key selection criteria}\\
Prefer keys that are:
\begin{itemize}[leftmargin=*,nosep]
  \item stable over time,
  \item always available at creation time,
  \item minimal (no unnecessary attributes),
  \item and meaningful only if the meaning is truly stable.
\end{itemize}
\end{principlebox}

If a key can change, you risk cascading updates to every reference.
If a key is not available when the entity is created, you will be forced to invent placeholder values.
Both outcomes degrade correctness.

\begin{figure}[htbp]
\centering
\includegraphics[width=0.75\textwidth]{key-types.pdf}
\caption{Primary, alternate, and surrogate keys all point to the same entity identity.}
\end{figure}

\begin{notebox}
Do not use names as primary keys unless the domain explicitly treats them as stable identifiers.
Names change, collide, and create hidden duplicates.
\end{notebox}

\section{Uniqueness beyond the primary key}
Primary keys are not the only uniqueness rules.
Many domains have alternate identifiers that must be unique when present.
These are sometimes called alternate or candidate keys.

\subsection{Alternate and candidate keys}
A \term{candidate key} is any set of attributes that could serve as a primary key.
Only one is chosen as the primary key, but the others still carry identity meaning.

\subsection{Partial uniqueness}
Some identifiers are optional but must be unique when present.
This is common with email addresses, national IDs, or barcodes.
The rule is not “always present” but “never duplicated.”

\begin{examplebox}
\textbf{Unique-when-present}\par\smallskip
\[
\rel{Member}(\underline{member\_id}, name, email? \;[\text{unique}])
\]
Email is optional, but if two members have the same email, the identity model is broken.
\end{examplebox}

\section{Composite keys and association relations}
Composite keys are common in association relations and in history tables.
They combine two or more attributes to form identity.

The most important example is a many-to-many association.
If a student can enroll in many courses, and a course can have many students, then the enrollment is identified by the pair.

\begin{examplebox}
\textbf{University enrollment identity}\par\smallskip
\[
\begin{aligned}
\rel{Enrollment}(&\underline{student\_id} \rightarrow \rel{Student}, \\
&\underline{course\_id} \rightarrow \rel{Course}, \\
&enrolled\_at, grade?)
\end{aligned}
\]
The pair (student\_id, course\_id) is the identity of the enrollment.
\end{examplebox}

Composite keys can also include a time attribute.
If you track historical states, the identity might be (entity\_id, effective\_from).
This keeps each historical record distinct.

\section{Key availability in workflows}
Key choice must fit how the system is used.
In a clinic, you may need to create a patient record before any external identifier exists.
In a library, a copy might be catalogued before its barcode is printed.
These are workflow facts, not technical details.

\begin{examplebox}
\textbf{Copy identity in a library}\\
If a barcode is assigned later, a surrogate key avoids placeholders:
\[
\rel{Copy}(\underline{copy\_id}, barcode \;[\text{unique}], status, title\_id \rightarrow \rel{Title})
\]
Barcode is unique when present, but not required for identity at creation time.
\end{examplebox}

The same logic applies to university enrollment.
The enrollment exists when the student registers, not when a final grade is assigned.
So the grade must not be part of the key.

\section{Choosing keys in clinic, library, and university}
Key decisions become clearer when you compare domains.

Clinic:
patients and doctors often need surrogate keys because external identifiers are inconsistent.
Appointments typically use a surrogate key because multiple appointments can share time or participants.

Library:
titles can sometimes use a natural key (ISBN), but copies almost never should.
Loans and reservations are events; a surrogate key or a composite key with a timestamp both work, depending on how the domain defines uniqueness.

University:
students often have an institutional id that is stable and can be a natural key.
Courses often have a code, but if codes are reused across years, a surrogate key is safer.
Enrollment is usually a composite key (student\_id, course\_id, term) if the domain allows repeat enrollment across terms.

\section{Worked mini-example: university keys}
A university model forces you to make identity choices early.
Students often have a stable institutional id, but courses may have reused codes.
Enrollment identity depends on the rule for retaking courses.

\begin{examplebox}
\textbf{University keys discussion}\par\smallskip
\[
\rel{Student}(\underline{student\_id}, name, email? \;[\text{unique}])
\]
\[
\rel{Course}(\underline{course\_id}, code \;[\text{unique}], title)
\]
\[
\begin{aligned}
\rel{Enrollment}(&\underline{student\_id} \rightarrow \rel{Student}, \\
&\underline{course\_id} \rightarrow \rel{Course}, \\
&\underline{term}, status)
\end{aligned}
\]
If the same course code can be reused across years, \texttt{course\_id} should be a surrogate key.
\end{examplebox}

\begin{examplebox}
\textbf{Enrollment with term}\\
If students can retake courses, include term in the identity:
\[
\rel{Enrollment}(\underline{student\_id} \rightarrow \rel{Student}, \underline{course\_id} \rightarrow \rel{Course}, \underline{term}, grade?)
\]
This preserves the rule “one enrollment per student, course, and term.”
\end{examplebox}

\section{Common mistakes}
Identity mistakes are easy to make and hard to fix later.

Mistake: Using email or name as a primary key.\\
Fix: Use a surrogate key and enforce email uniqueness when present.

Mistake: Using composite keys everywhere “for purity.”\\
Fix: Use composite keys only when the domain rule truly defines identity.

Mistake: Ignoring uniqueness of domain identifiers (barcode, national ID).\\
Fix: State alternate keys as explicit uniqueness constraints.

Mistake: Putting status or time attributes in the key.\\
Fix: Keep status and timestamps as non-key attributes unless the rule demands them.

\section{Mini checklist}
Before you finalize keys, check:
\begin{itemize}[leftmargin=*,nosep]
  \item Is every entity’s identity explicit and stable?
  \item Is the key available at creation time?
  \item If you use a natural key, is it truly unique by rule?
  \item Are alternate identifiers stated as uniqueness constraints?
  \item Do association relations have composite keys that match their uniqueness rules?
  \item Are time or status attributes excluded from keys unless the rule requires them?
  \item Can you explain how merges or splits would affect references?
  \item Are keys protected from changes in UI or workflow convenience?
\end{itemize}

\begin{chaptersummary}
Keys define identity and make relationships possible.
Natural keys can be powerful when stable and universal, but surrogate keys often provide safer identity.
Composite keys are common in association relations and history structures.
Key choices must match the domain rules and the workflow that creates records.
If identity is unclear, the rest of the schema cannot be trusted.
\end{chaptersummary}

\begin{exercisebox}
Answer in full sentences and then write the corresponding schema fragments.
\begin{enumerate}[label=\arabic*.]
  \item For a clinic, propose natural identifiers for Patient and Doctor and explain why they are or are not suitable as keys.
  \item For a library, decide whether Title should use ISBN as a key and justify your choice.
  \item Design a composite key for an association relation in a university domain and state the rule it enforces.
  \item Give one example where a key attribute might change and explain the risk.
  \item Propose a surrogate key strategy for a domain you know and list two benefits.
  \item Show a relation where term or time belongs in the key, and explain the uniqueness rule.
  \item Identify one attribute that should never be part of a key and justify why.
  \item Write a uniqueness rule that is “unique when present” and show it in schema notation.
  \item Describe a merge scenario (duplicates) and explain how your key choice makes it manageable.
\end{enumerate}
\end{exercisebox}

In the next chapter we translate conceptual models into relational schemas.
The mapping rules build directly on the identity and uniqueness choices made here.
